\chapter*{Lista de abreviaturas e siglas}

% Gerada pelo pacote acronym: no texto, use \ac{X} (1º uso expande),
% \acs{X} (só a sigla), \acl{X} (só o nome) e \acp{X} (plural).
% Com a opção printonlyused, só aparecem as siglas usadas no texto.
% \acroextra{...} aparece apenas nesta lista.
\begin{acronym}[QSPI]
  \acro{ADC}{conversor analógico-digital\acroextra{ (\emph{Analog-to-Digital Converter})}}
  \acro{AHB}{\emph{Advanced High-performance Bus}}
  \acro{API}{interface de programação de aplicações\acroextra{ (\emph{Application Programming Interface})}}
  \acroplural{API}[APIs]{interfaces de programação de aplicações}
  \acro{APB}{\emph{Advanced Peripheral Bus}}
  \acro{ASIC}{circuito integrado de aplicação específica\acroextra{ (\emph{Application-Specific Integrated Circuit})}}
  \acroplural{ASIC}[ASICs]{circuitos integrados de aplicação específica}
  \acro{BSS}{\emph{Block Started by Symbol}\acroextra{ (seção de dados não inicializados)}}
  \acro{CAN}{\emph{Controller Area Network}}
  \acro{CDC}{\emph{Communications Device Class}\acroextra{ (classe de dispositivo USB)}}
  \acro{CPU}{unidade central de processamento\acroextra{ (\emph{Central Processing Unit})}}
  \acro{DAQ}{dispositivo de aquisição de dados\acroextra{ (\emph{Data Acquisition})}}
  \acro{DMA}{acesso direto à memória\acroextra{ (\emph{Direct Memory Access})}}
  \acro{DREQ}{sinal de requisição de dados\acroextra{ (\emph{Data Request})}}
  \acro{DWT}{\emph{Data Watchpoint and Trace}\acroextra{ (unidade do Cortex-M com contador de ciclos)}}
  \acro{ELF}{\emph{Executable and Linkable Format}}
  \acro{FIFO}{fila \emph{first-in, first-out}}
  \acroplural{FIFO}[FIFOs]{filas \emph{first-in, first-out}}
  \acro{FPGA}{\emph{Field-Programmable Gate Array}}
  \acroplural{FPGA}[FPGAs]{\emph{Field-Programmable Gate Arrays}}
  \acro{GPIO}{entrada e saída de propósito geral\acroextra{ (\emph{General-Purpose Input/Output})}}
  \acroplural{GPIO}[GPIOs]{entradas e saídas de propósito geral}
  \acro{I2C}[I2C]{\emph{Inter-Integrated Circuit}}
  \acro{IC}{intervalo de confiança}
  \acro{IoT}{Internet das Coisas\acroextra{ (\emph{Internet of Things})}}
  \acro{IRQ}{requisição de interrupção\acroextra{ (\emph{Interrupt Request})}}
  \acroplural{IRQ}[IRQs]{requisições de interrupção}
  \acro{ISR}{rotina de serviço de interrupção\acroextra{ (\emph{Interrupt Service Routine})}}
  \acroplural{ISR}[ISRs]{rotinas de serviço de interrupção}
  \acro{MCU}{microcontrolador\acroextra{ (\emph{Microcontroller Unit})}}
  \acroplural{MCU}[MCUs]{microcontroladores}
  \acro{PIO}{entrada/saída programável\acroextra{ (\emph{Programmable Input/Output})}}
  \acro{QSPI}{\emph{Quad Serial Peripheral Interface}}
  \acro{RLE}{codificação por comprimento de sequência\acroextra{ (\emph{Run-Length Encoding})}}
  \acro{RTOS}{sistema operacional de tempo real\acroextra{ (\emph{Real-Time Operating System})}}
  \acro{SDK}{\emph{Software Development Kit}}
  \acro{SIO}{bloco de E/S de ciclo único\acroextra{ (\emph{Single-cycle I/O})}}
  \acro{SM}{máquina de estados do PIO\acroextra{ (\emph{State Machine})}}
  \acroplural{SM}[SMs]{máquinas de estados}
  \acro{SPI}{\emph{Serial Peripheral Interface}}
  \acro{SRAM}{memória estática de acesso aleatório\acroextra{ (\emph{Static Random-Access Memory})}}
  \acro{UART}{\emph{Universal Asynchronous Receiver-Transmitter}}
  \acro{USB}{\emph{Universal Serial Bus}}
  \acro{XIP}{execução direta da memória flash\acroextra{ (\emph{Execute-In-Place})}}
\end{acronym}
